% Template for Elsevier CRC journal article
% version 1.2 dated 09 May 2011

% This file (c) 2009-2011 Elsevier Ltd.  Modifications may be freely made,
% provided the edited file is saved under a different name

% This file contains modifications for Transportation Research Procedia

% Changes since version 1.1
% - added "procedia" option compliant with ecrc.sty version 1.2a
%   (makes the layout approximately the same as the Word CRC template)
% - added example for generating copyright line in abstract

%-----------------------------------------------------------------------------------

%% This template uses the elsarticle.cls document class and the extension package ecrc.sty
%% For full documentation on usage of elsarticle.cls, consult the documentation "elsdoc.pdf"
%% Further resources available at http://www.elsevier.com/latex

%-----------------------------------------------------------------------------------

%%%%%%%%%%%%%%%%%%%%%%%%%%%%%%%%%%%%%%%%%%%%%%%%%%%%%%%%%%%%%%
%%%%%%%%%%%%%%%%%%%%%%%%%%%%%%%%%%%%%%%%%%%%%%%%%%%%%%%%%%%%%%
%%                                                          %%
%% Important note on usage                                  %%
%% -----------------------                                  %%
%% This file should normally be compiled with PDFLaTeX      %%
%% Using standard LaTeX should work but may produce clashes %%
%%                                                          %%
%%%%%%%%%%%%%%%%%%%%%%%%%%%%%%%%%%%%%%%%%%%%%%%%%%%%%%%%%%%%%%
%%%%%%%%%%%%%%%%%%%%%%%%%%%%%%%%%%%%%%%%%%%%%%%%%%%%%%%%%%%%%%

%% The '3p' and 'times' class options of elsarticle are used for Elsevier CRC
%% The 'procedia' option causes ecrc to approximate to the Word template
\documentclass[5p,times,procedia]{elsarticle}
\flushbottom

%% The `ecrc' package must be called to make the CRC functionality available
\usepackage{ecrc}
%\usepackage{amsmath}
% tighten float spacing (saves surprisingly much)
\setlength{\textfloatsep}{8pt plus 2pt minus 2pt} % space between top/bot floats and text
\setlength{\intextsep}{8pt plus 2pt minus 2pt}   % space for in-text floats
\setlength{\floatsep}{6pt plus 2pt minus 2pt}    % space between floats
\setlength{\abovecaptionskip}{3pt}
\setlength{\belowcaptionskip}{0pt}

%% The ecrc package defines commands needed for running heads and logos.
%% For running heads, you can set the journal name, the volume, the starting page and the authors

%% set the volume if you know. Otherwise `00'
\volume{00}

%% set the starting page if not 1
\firstpage{1}

%% Give the name of the journal
\journalname{Procedia CIRP}

%% Give the author list to appear in the running head
%% Example \runauth{C.V. Radhakrishnan et al.}
\runauth{Author name}

%% The choice of journal logo is determined by the \jid and \jnltitlelogo commands.
%% A user-supplied logo with the name <\jid>logo.pdf will be inserted if present.
%% e.g. if \jid{yspmi} the system will look for a file yspmilogo.pdf
%% Otherwise the content of \jnltitlelogo will be set between horizontal lines as a default logo

%% Give the abbreviation of the Journal.
\jid{trpro}

%% Give a short journal name for the dummy logo (if needed)
%\jnltitlelogo{Transportation Research}

%% Hereafter the template follows `elsarticle'.
%% For more details see the existing template files elsarticle-template-harv.tex and elsarticle-template-num.tex.

%% Elsevier CRC generally uses a numbered reference style
%% For this, the conventions of elsarticle-template-num.tex should be followed (included below)
%% If using BibTeX, use the style file elsarticle-num.bst

%% End of ecrc-specific commands
%%%%%%%%%%%%%%%%%%%%%%%%%%%%%%%%%%%%%%%%%%%%%%%%%%%%%%%%%%%%%%%%%%%%%%%%%%
\usepackage{siunitx}
%% The amssymb package provides various useful mathematical symbols

\usepackage{amssymb}
%% The amsthm package provides extended theorem environments
%% \usepackage{amsthm}

%% The lineno packages adds line numbers. Start line numbering with
%% \begin{linenumbers}, end it with \end{linenumbers}. Or switch it on
%% for the whole article with \linenumbers after \end{frontmatter}.
%% \usepackage{lineno}

%% natbib.sty is loaded by default. However, natbib options can be
%% provided with \biboptions{...} command. Following options are
%% valid:

%%   round  -  round parentheses are used (default)
%%   square -  square brackets are used   [option]
%%   curly  -  curly braces are used      {option}
%%   angle  -  angle brackets are used    <option>
%%   semicolon  -  multiple citations separated by semi-colon
%%   colon  - same as semicolon, an earlier confusion
%%   comma  -  separated by comma
%%   numbers-  selects numerical citations
%%   super  -  numerical citations as superscripts
%%   sort   -  sorts multiple citations according to order in ref. list
%%   sort&compress   -  like sort, but also compresses numerical citations
%%   compress - compresses without sorting
%%
%\biboptions{authoryear}

% \biboptions{}

% if you have landscape tables
\usepackage[figuresright]{rotating}
%\usepackage{harvard}
% put your own definitions here:x
%   \newcommand{\cZ}{\cal{Z}}
%   \newtheorem{def}{Definition}[section]
%   ...

% add words to TeX's hyphenation exception list
%\hyphenation{author another created financial paper re-commend-ed Post-Script}

% declarations for front matter

\usepackage[bookmarks=false]{hyperref}
    \hypersetup{colorlinks,
      linkcolor=blue,
      citecolor=blue,
      urlcolor=blue}


\begin{document}
\begin{frontmatter}

%% Title, authors and addresses

%% use the tnoteref command within \title for footnotes;
%% use the tnotetext command for the associated footnote;
%% use the fnref command within \author or \address for footnotes;
%% use the fntext command for the associated footnote;
%% use the corref command within \author for corresponding author footnotes;
%% use the cortext command for the associated footnote;
%% use the ead command for the email address,
%% and the form \ead[url] for the home page:
%%
%% \title{Title\tnoteref{label1}}
%% \tnotetext[label1]{}
%% \author{Name\corref{cor1}\fnref{label2}}
%% \ead{email address}
%% \ead[url]{home page}
%% \fntext[label2]{}
%% \cortext[cor1]{}
%% \address{Address\fnref{label3}}
%% \fntext[label3]{}

\dochead{10th Conference on High Performance Cutting (CIRP-HPC 2026)}%

\title{Prediction of Wire Breakage in Wire Electrical Discharge Machining by using a Thermomechanical Damage Model}

%% use optional labels to link authors explicitly to addresses:
%% \author[label1,label2]{<author name>}
%% \address[label1]{<address>}
%% \address[label2]{<address>}


\author[a,b]{Eduardo González-Sánchez\corref{cor1}}
\author[a,c]{Nicolás Muñoz} 
\author[a]{Konrad Wegener}
\author[a]{Davide Beretta}
\author[a]{Amir Malakizadi}

%\fntext[equal]{These authors contributed equally to this work.}
%\ead{author@institute.xxx}

\address[a]{inspire AG, Technoparkstrasse 1, Zurich, 8005, Switzerland}
\address[b]{Chair of Artificial Intelligence in Mechanics and Manufacturing, ETH Zurich, Rämistrasse 101, Zurich, 8092, Switzerland}
\address[c]{Tecnun School of Engineering, Universidad de Navarra, Manuel Lardizabal 13, San Sebastián, 20018,  Spain}

\cortext[cor1]{* Corresponding author. {\it E-mail address:} eduardo.gonzalez@inspire.ch}


\begin{abstract}
%% Text of abstract
Wire breakage is the main limitation to productivity in Wire Electrical Discharge Machining (WEDM). In this process, the wire is continuously fed at constant tension while undergoing repeated heating from electrical discharges, which can lead to wire breakage and interrupt machining. This study examines the thermomechanical causes of wire failure, focusing on temperature-driven softening and time-dependent strength reduction. Brass wires of various diameters were tested under different mechanical and thermal loads. The temperature of the wire was controlled by closed-loop Joule heating and tracked with infrared thermography. Microstructural analysis via Scanning Electron Microscopy (SEM) confirms that the wire can undergo dynamic recrystallization within the time and temperature ranges relevant to the rough-cut WEDM process. Experimental results indicate that wire breakage is governed by a thermally activated loss of load-carrying capacity that depends on temperature, stress, and time. Based on these findings, this study introduces a damage-accumulation model with power-law stress and Arrhenius-type temperature dependency, calibrated using experimental data obtained under varying stress levels and heating rates. The model is then implemented in the SPARC finite-element simulation framework for WEDM, which uses a computed transient thermal field to evolve a segment-wise damage variable and predict wire breakage. Preliminary simulations show qualitative agreement with observed machine behavior, distinguishing between safe and critical operating conditions.
\end{abstract}

\begin{keyword}
edm \sep wire \sep brass \sep process \sep model \sep drx
%% keywords here, in the form: keyword \sep keyword

%% PACS codes here, in the form: \PACS code \sep code

%% MSC codes here, in the form: \MSC code \sep code
%% or \MSC[2008] code \sep code (2000 is the default)

\end{keyword}
%\cortext[cor1]{Corresponding author. Tel.: +0-000-000-0000 ; fax: +0-000-000-0000.}

\end{frontmatter}

%%
%% Start line numbering here if you want
%%
% \linenumbers

%% main text

%\enlargethispage{-7mm}
\section{Introduction}
\label{introduction}

Wire Electrical Discharge Machining (WEDM) is a machining process widely used for producing components with complex geometries and high precision ($1\sim10 ~\mathrm{\mu m}$) of electrically conductive materials. In this process, material removal occurs through a series of controlled electrical discharges between a continuously fed wire (the tool electrode) and the workpiece, enabling the machining of hard materials without direct mechanical contact. Due to its precision, WEDM has become an essential technology in high-value sectors such as aerospace, biomedical, automotive, and precision tooling \cite{slatineanu2020,ho2004state}. Despite significant advances in recent years, wire breakage remains one of the most critical limitations affecting process efficiency: when failure occurs, machining must be interrupted and restarted, leading to substantial time losses and increased production costs. The mechanisms responsible for wire failure have been the focus of extensive research. Some investigations suggest that the combination of spark pressure and wire tension drives the propagation of erosion-induced cracks along the wire, thereby reducing its effective cross-sectional area and ultimately leading to failure \cite{wang2023,luo1999}. However, thermally induced ductile fracture (\autoref{gr1.jpg}), resulting from loss of tensile-strength caused by excessive heat accumulation, is by far the most common failure mode in brass or copper wires \cite{schacht2004}.

\begin{figure}[t]
    \centering
    \includegraphics[width=0.8\columnwidth]{gr1.jpg}
    \caption{Necking of a brass wire (CuZn37 LP900, $\varnothing 0.2~\mathrm{mm}$) after ductile fracture during WEDM of a $50~\mathrm{mm}$ steel workpiece on an Agie Charmilles CUT E 350.}
    \label{gr1.jpg}
\end{figure}


Simple temperature and/or spark energy threshold criteria are insufficient for predicting wire failure, as breakage depends on the cumulative thermomechanical history rather than instantaneous conditions alone. For instance, a concentration of discharges in a localized region can lead to excessive local heating and rapid strength degradation, even if the average power remains within nominal limits. The importance of spark location tracking has been recognized in industrial practice, with approaches using fast FPGAs to monitor spark locations in real-time and suppress potentially dangerous discharges \cite{di_campli_real-time_2020,damario2019patent}. Nevertheless, the space of possible control strategies remains largely unexplored, with potential for sophisticated approaches that dynamically adjust process parameters based on the wire's accumulated thermal and mechanical history. A damage accumulation model is therefore critical to enable the development and testing of advanced control strategies. 

This work builds on the SPARC simulation platform \cite{sparc}, a modular finite-element framework designed to simulate the main cut WEDM process for process optimization, including the transient thermal field of the wire electrode. The model discretizes the wire into segments and solves the 1D heat equation, accounting for conduction, convection, and source terms. \autoref{gr2} illustrates the thermal problem as simulated by SPARC. The discharge current $I$ comes from two components, $I_1$ and $I_2$, flowing from the upper and lower contacts respectively, each depending on the path resistance, which in turn depends on the spark location and the temperature-dependent resistivity. Although the thermal model currently implemented in SPARC delivers detailed temperature profiles, the criterion for wire breakage relies exclusively on a simple temperature threshold. The present study aims to develop and calibrate a physics-based cumulative damage model capable of predicting wire breakage under realistic loading conditions. Laboratory tests were performed to replicate the mechanical tension and thermal conditions experienced by the wire during machining, allowing the measurement of temperature, stress, and time to failure. Based on these results, this work presents a damage rate model $\dot{D}(T, \sigma)$ enabling accurate prediction of wire breakage based on the wire's cumulative thermomechanical history.

\begin{figure}[t]
    \centering
    \includegraphics[height=5cm]{gr2a.jpg}
    \includegraphics[height=5cm]{gr2b.png}
    \caption{Left: Schematic of the wire's current distribution during a discharge. Right: Simulated thermal evolution of a wire segment as it moves through the cutting zone obtained from the SPARC simulation.}
    \label{gr2}
\end{figure}


\section{Experimental Methods and Damage Model Calibration}
\label{experimental}
\subsection{Measurement Setup}
\label{setup}
An experimental setup was designed to measure the ultimate tensile strength of a electrode wire under the typical thermomechanical WEDM operating conditions. Specifically, the system applies a uniaxial tensile load to the wire while simultaneously heating it by Joule effect with a direct current, and records the wire temperature as a function of time until fracture occurs. \autoref{gr4} shows a schematic of the system, including an infrared thermal camera (X6801sc, \emph{Teledyne Technologies}), a microcontroller (\emph{Arduino}), a power supply (6225A, \emph{PeakTech}), a MOSFET (IRLB3034PBF, \emph{Infineon Technologies}), a computer, and the wire under test. All control and data acquisition processes are handled by a Python interface, including real-time monitoring and storage of the following experimental variables: time and wire's maximum temperature. The infrared camera records the wire temperature up to $400~\mathrm{^\circ C}$ at a sampling frequency of $520~\mathrm{Hz}$ (\autoref{gr5}), while the wire's maximum temperature is determined by identifying the highest grayscale value (in the range $0-255$) among the camera pixels and linearly interpolating it over the selected temperature range. This temperature serves as the feedback signal for the PID controller. The Python interface calculates the error between the setpoint and the measured temperature and generates a digital Pulse Width Modulation (PWM) signal for the micro-controller. The micro-controller converts this digital signal into an analog PWM output at a frequency of $1~\mathrm{kHz}$ to modulate the MOSFET gate voltage, which in turn regulates the current supplied by the power generator to the wire in the range $0- 5~\mathrm{A}$, corresponding to a maximum peak current density of ca. $10^{2}~\mathrm{A/mm^2}$. To minimize surface reflections and accurately capture the wire temperature, the wire surface was manually coated with matte black paint (HERP-HT-MWIR-BK-11, \emph{Thermal Focus}), resulting in an uncertainty of $\pm 10~\mathrm{^\circ C}$ for the wire's maximum temperature. It is worth noting that the thermal camera software (FLIR ResearchIR) used in this study imposed a limitation on the selectable temperature range, requiring the user to choose either the low range ($20–200~\mathrm{^\circ C}$) or the high range ($150–400~\mathrm{^\circ C}$). Since wire damage is negligible below $150~\mathrm{^\circ C}$ under the relevant timescales of this process, all recordings used the high temperature range.


\begin{figure}[t] % 
\centering 
\includegraphics[width=1.0\columnwidth]{gr4.png} 
\caption{Schematic diagram of the experimental setup.}
\label{gr4}
\end{figure}

\begin{figure}[t] 
\centering
\includegraphics[width=1.0\columnwidth]{gr5.png}
\caption{Left: Photo of the experimental setup. Right: Thermographic image of a representative wire specimen, overlaid on a grayscale image of the experimental setup.}
\label{gr5}
\end{figure}

\subsection{Materials and Sample Preparation}
This work focuses on commercial cold-drawn brass wires with a nominal tensile strength of $900~\mathrm{MPa}$ (CuZn37 LP900, \emph{Thermocompact}) in two diameters, $0.15~\mathrm{mm}$ and $0.20~\mathrm{mm}$, used as received. Samples for Scanning Electron Microscopy (Zeiss EVO 10) were prepared according to the following procedure. First, the samples were embedded in a resin (QPREP UV 55, \emph{ATM}) and then treated in vacuum to eliminate entrapped air bubbles. The resin was then cured under UV light for $15 ~ \mathrm{min}$ to obtain a solid embedding. The exposed surface was ground using SiC abrasive papers with progressively finer grain sizes of $12 ~ \mathrm{\mu m}$, $9~\mathrm{\mu m}$, and $5~\mathrm{\mu m}$, at $300~\mathrm{rpm}$. Subsequent polishing was performed using a $1 ~ \mathrm{\mu m}$ abrasive diamond disc at $150~\mathrm{rpm}$ to achieve a mirror-like finish. A final vibro-polishing step was carried out for approximately $3 ~ \mathrm{h}$ using a Cu–Zn suspension to remove residual micro-scratches. The microstructure was finally revealed by chemical etching with Copper(I) chloride (\emph{ATM}), followed by the deposition of a thin conductive carbon coating to suppress surface charging during SEM imaging.


\subsection{Isothermal Annealing Test}
To evaluate the effect of temperature on wire strength, $16$ series of isothermal annealing tests were conducted at six different temperatures, all below the $\alpha\rightarrow\beta$ phase transition \cite{asm2008}. Samples were heated at a constant temperature, then cooled to room temperature, and subsequently subjected to a monotonically increasing tensile load until failure to determine the ultimate tensile strength. 

\subsection{Electromigration Effect Tests}

One concern is whether current induced electromigration accelerates failure for a similar thermal and tensile history. To investigate this, a comparative study between two wire heating methods was conducted: (i) Joule heating by direct current, with current densities in the range $50$--$300~\mathrm{A/mm^2}$, and (ii) convective heating using a heat gun pointed directly at the wire while recording its temperature with the infrared camera. A total of 8 experiments were conducted, applying a constant uniaxial tensile load of $320~\mathrm{MPa}$ during the tests. 

It should be noted that the peak current densities in a typical WEDM process may be one to two orders of magnitude higher than that applied with this work's experimental setup. Nevertheless, the time-averaged current density remains comparable. Therefore, it is to assume that electromigration is also negligible in a typical WEDM machine for current densities smaller than ca. $10^4~\mathrm{A/mm^2}$ \cite{ho1989}.

\subsection{Thermomechanical Breaking Tests}

The effect of the samples' temperature history was investigated through 48 experiments conducted under constant uniaxial load (16 tests for each stress level: $320$, $450$, and $580~\mathrm{MPa}$). For each test, linear DC current rates from $0.25~\mathrm{A/s}$ to $50~\mathrm{A/s}$  were applied, resulting in heating rates ranging from $4~\mathrm{^\circ C/s}$ to $400~\mathrm{^\circ C/s}$ until failure occurred. Throughout each trial, time and temperature evolution were continuously recorded using the setup described in \autoref{setup}.


\subsection{Damage Model and Data Fitting}\label{model}

Wire failure in WEDM is governed by a progressive, thermally activated loss of tensile strength. SEM observations (\autoref{gr7}) confirm that dynamic recrystallization may occur within the time and temperature ranges of the experiments. While recovery is not directly observable via SEM imaging, the evidence of recrystallization suggests that recovery has already taken place.~\cite{humphreys2017}.

A scalar damage variable $D$ is introduced to represent the cumulative fractional loss of load-carrying capacity, with $D=0$ for the pristine, undamaged wire and $D=1$ for a fully damaged wire at rupture. Since tracking strain is not feasible in the current experimental setup nor in the SPARC simulation framework, damage evolution is formulated without explicit strain dependence, relying instead on stress and temperature history to predict rupture. The damage rate is:

\begin{equation}\label{eq:damage_rate}
\dot{D} = k\, \sigma^{n} \exp\left(-\frac{Q}{RT}\right)
\end{equation}
where $k$ is a rate constant, $\sigma$ the applied stress, $n$ the stress exponent, $Q$ the activation energy, $R$ the gas constant, and $T$ the absolute temperature. The power-law stress dependence $\sigma^n$ is standard in rupture-life correlations~\cite{zhang2010creep}, while the Arrhenius term $\exp(-Q/RT)$ captures the thermally activated nature of the underlying restoration processes, the same functional form used in the derivation of the Larson--Miller parameter~\cite{larsonmiller1952}. This formulation corresponds to a separable evolution law with no explicit dependence on $D$, which Stigh~\cite{stigh2006} showed to be mathematically equivalent to Robinson's linear life-fraction rule ~\cite{robinson1952}.

The Kachanov effective-stress formulation $(\sigma/(1-D))^{n}$~\cite{kachanov} was tested but found to degrade predictive accuracy. This term models stress concentration on the remaining load-bearing cross-section as damage accumulates, accelerating failure. In the present experiments, however, this final instability manifests as necking occurring within a very short time, i.e. ca. a single camera frame ($<5~\mathrm{ms}$), which is negligible compared to the $0.5$--$10~\mathrm{s}$ heating duration. The linear formulation is therefore appropriate.

The model parameters $(k, n, Q)$ are fitted to a dataset of 48 experiments by minimizing:
\begin{equation}
C = \frac{1}{N}\sum_{i=1}^{N} w_i \left(D_{\mathrm{final},i} - 1\right)^2
\end{equation}
where $D_{\mathrm{final},i}$ is the accumulated damage at the experimental rupture time $t^{*}_{\mathrm{exp},i}$ for test $i$, and $w_i = 1/t^{*}_{\mathrm{exp},i}$ is a weighting factor. The inverse-time weighting is necessary because $D_{\mathrm{final}} = \int_0^{t^{*}_{\mathrm{exp}}} \dot{D}\,dt$. Therefore, a given error in the damage rate $\dot{D}$ accumulates over the integration interval, producing a residual that scales with $t^{*}_{\mathrm{exp}}$. Without the weighting factor $w_i$, long-duration tests would dominate the cost function, causing the optimizer to neglect short-duration failures. Optimization was performed using the Differential Evolution algorithm via \texttt{scipy.optimize}, with forward Euler integration at a $2~\mathrm{ms}$ time step.


\section{Results and Discussion}
\label{results}

\subsection{Isothermal Annealing}
The results in \autoref{gr6} show that higher annealing temperatures lead to lower residual tensile strength, and the strength reduction is only ca. 10$\%$ for annealing temperatures below $150~\mathrm{^\circ C}$ and annealing times up to one minute. This is compatible with dynamic recrystallization of heavily cold worked brass such as the one subject of this study. Besides, fracture consistently occurs at the hottest region (i.e., the wire midpoint), confirming that localized thermal damage governs the failure location. SEM microstructural analysis of samples before and after annealing (\autoref{gr7}) reveals that exposure to $350~\mathrm{^\circ C}$ for just $3~\mathrm{s}$ is sufficient to initiate recrystallization, resulting in a tensile-strength reduction of approximately 60\% as shown in \autoref{gr6}.
\begin{figure}[t] 
\centering 
\includegraphics[width=0.8\columnwidth]{gr6.png}  
\caption{Tensile strength of the brass wires at room temperature after annealing at different temperatures and times.}
\label{gr6}
\end{figure}

\begin{figure}[t]
\centering
\includegraphics[width=0.3\textwidth]{gr7a.png}
\hspace{0.1cm}
\includegraphics[width=0.3\textwidth]{gr7b.png}
\caption{Representative SEM images of wire cross sections. Top: pristine wire. Bottom: wire heated in the experimental setup at $350\,^\circ\mathrm{C}$ for $3~\mathrm{s}$.}
\label{gr7}
\end{figure}

\subsection{Electromigration Effects}
A potential concern in using Joule heating to study wire failure is whether electromigration—the transport of metal atoms by electron momentum transfer—could accelerate recrystallization or otherwise affect the wire's mechanical integrity independently of thermal effects \cite{wang2001}. As shown in \autoref{gr8}, the failure temperatures obtained by both methods are comparable (within 5$\%$) for applied current densities in the range $50$--$300~\mathrm{A/mm^2}$. This result indicates that electromigration effects are negligible under the considered experimental conditions, and that the observed thermal degradation is primarily governed by temperature. 

\begin{figure}[t] 
\centering 
\includegraphics[width=0.8\columnwidth]{gr8.png}  
\caption{Comparison of the temperature evolution up to wire breaking for representative samples under Joule and convection heating. Time zero corresponds to the moment when the wire's highest temperature first exceeds $150~\mathrm{^\circ C}$.}
\label{gr8}
\end{figure}

\subsection{Thermomechanical Breaking and Model Calibration}
The experimental results, presented in \autoref{gr9}, reveal that the heating rate significantly affects the breaking temperature: higher heating rates lead to higher breaking temperatures. This behavior is consistent with the time-dependent nature of the damage accumulation mechanism modeled in this work. At lower heating rates, the wire spends more time at elevated temperatures, allowing greater cumulative damage to accrue before reaching the critical breakage threshold ($D=1$). The damage model parameters were calibrated on the complete dataset of 48 thermomechanical breaking tests using the weighted least-squares objective defined in \autoref{model}. The resulting best-fit parameters are reported in \autoref{Tab1}, corresponding to a stress exponent of $n=6.73$ and an activation energy of $Q=143.1~\mathrm{kJ/mol}$, which are consistent with the values found in the literature \cite{pernis2010,bursikova2002}. \autoref{gr9} shows that the calibrated model captures the observed dependence of rupture temperature on heating rate across all three stress levels. For completeness, the agreement on the calibration dataset is quantified by comparing the experimental rupture times $t^{*}_{\mathrm{exp}}$ to the model-predicted rupture times $t^{*}_{\mathrm{pred}}$, defined as the first time at which the integrated damage reaches $D(t)=1$. In cases where the damage had not reached unity at the experimental rupture time (i.e., $D(t^{*}_{\mathrm{exp}})<1$), the measured temperature history was continued by a linear extrapolation using the average heating rate, and the damage integration was extended until $D=1$ in order to obtain a corresponding $t^{*}_{\mathrm{pred}}$. Using this definition, the model yields a mean absolute percentage error of 10.9\% and a root mean square error of $1.38~\mathrm{s}$ over the 48 tests. This performance is considered sufficient for the purpose of this work, which is to capture a representative damage accumulation mechanism for wire breakage in WEDM, enabling the testing of advanced control strategies within the SPARC simulation framework.

\begin{figure}[t] 
\centering 
\includegraphics[width=0.8\columnwidth]{gr9.png}  
\caption{Temperature histories and rupture points for 48 thermomechanical tests at constant tensile stress (\(\sigma = 320, 450,\) and \(580~\mathrm{MPa}\)). Thin curves show the measured maximum wire temperature \(T_{\max}(t)\) for each test; symbols mark the rupture point \((\Delta t, T^*)\), where \(\Delta t = t^* - t_{150}\) and \(t_{150}\) is the time when \(T_{\max}\) first exceeds \(150~^\circ\mathrm{C}\). Solid curves indicate the model-predicted rupture point (\(D=1\)) for each stress level using the fitted parameters in \autoref{Tab1}, computed assuming constant heating rates spanning the experimental range.
} 
\label{gr9}
\end{figure}

\begin{table}[h]
\caption{Damage model's parameters}
\begin{tabular*}{\hsize}{@{\extracolsep{\fill}}lll@{}}
\toprule 
Parameter & Units & Value \\
\colrule
k  & $s^{-1}\mathrm{MPa}^{-n}$ & 7.168\times 10^{-5}\\
Q  & $\mathrm{J/mol}$ & 143.1\times 10^{3}\\
n  & - & 6.734\\
\label{Tab1}
\botrule
\end{tabular*}
\end{table}


\subsection{Simulation of Wire Breakage in WEDM}
\begin{figure*}[t] 
\centering
\includegraphics[height=0.35\textwidth]{gr10a.png}
\hspace{0.05\textwidth}
\includegraphics[height=0.35\textwidth]{gr10b.png}
\caption{Temporal evolution of wire-segment temperature and cumulative damage predicted by the SPARC simulation for two discharge energy settings. The colored curve shows the segment temperature \(T(t)\), and the green curve shows the accumulated damage \(D(t)\) computed from the calibrated damage model. Time is measured from the instant the segment passes the upper power contact; the right axis reports the corresponding axial position relative to the upper contact based on the wire feed speed used in the simulation. Left: setting I15 (no rupture, \(D<1\)). Right: setting I17 (rupture when \(D=1\), marked by the cross).}



\label{gr10}
\end{figure*}

The calibrated damage model was implemented in the SPARC simulation framework to predict wire breakage under realistic WEDM conditions. \autoref{gr10} shows two representative simulation results illustrating the temporal evolution of temperature and cumulative damage for a wire segment as it traverses the cutting zone. The simulations were performed using the properties of a $0.20~\mathrm{mm}$ brass wire under two distinct discharge energy settings, denoted as I15 and I17, corresponding to standard parameters on an Agie Charmilles CUT E 350 machine. The I17 setting generates discharge pulses with higher peak current compared to the I15 setting. In the left panel (I15), the thermal loading remains moderate, and the damage variable $D$ increases gradually but stabilizes below the critical threshold, indicating that the wire survives. In contrast, the right panel (I17) depicts a scenario with higher energy input, where the temperature spikes are significantly higher, causing the damage to reach $D=1$, corresponding to wire rupture. These simulation predictions are consistent with experimental observations on the machine, where the wire survives under the I15 setting but breaks under the I17 settings. These initial results demonstrate that the model correctly distinguishes between safe and critical operating conditions. Further validation across a broader range of cutting conditions is planned to assess the model's general applicability in WEDM.




\section{Conclusions}
\label{conclusions}

This study presents a thermomechanical damage model for predicting wire breakage in Wire Electrical Discharge Machining. The main findings are:

\begin{itemize}
    \item Wire failure in WEDM is governed by temperature- and time-dependent strength degradation rather than instantaneous thermal thresholds. SEM analysis confirms that recrystallization initiates at temperatures as low as $350~\mathrm{^\circ C}$ within 3 seconds, leading to approximately 60\% reduction in tensile strength. 
    
    \item Electromigration effects are negligible under the studied conditions, as comparable failure temperatures were observed for both Joule heating and convective heating methods at current densities up to $300~\mathrm{A/mm^2}$.
    
    \item The heating rate significantly influences the breaking temperature: faster heating rates result in higher breaking temperatures, while slower rates allow more cumulative damage accumulation at lower temperatures.
    
    \item A damage accumulation model with power-law stress dependence and Arrhenius temperature dependence was successfully calibrated using 48 experiments under varying stress levels ($320$--$580~\mathrm{MPa}$) and heating rates ($4$--$400~\mathrm{^\circ C/s}$). The model parameters ($k$, $n$, $Q$) enable prediction of wire failure based on the complete thermal history.
    
\end{itemize}

The proposed damage model provides a foundation for future developments in WEDM process control. Validation under industrial cutting conditions is needed to confirm its applicability in real machining scenarios. However, the current model is sufficiently representative for the development of advanced control strategies that dynamically react based on accumulated damage, potentially enabling higher material removal rates while maintaining process stability. This work thus represents a significant step towards more intelligent and efficient WEDM operations.


\section*{Author contributions: CRediT}
\textbf{Nicolás Muñoz}:  Investigation, Writing - original draft. \textbf{Eduardo González-Sánchez}: Conceptualization, Investigation, Writing - original draft, Supervision. \textbf{Konrad Wegener}: Funding acquisition, Supervision, Writing - review and editing. \textbf{Davide Beretta}: Investigation, Project Administration, Supervision, Writing - review and editing. \textbf{Amir Malakizadi}: Writing - review and editing.


 \begin{thebibliography}{00}

%% \bibitem must have the following form:
%%   \bibitem{key}...
%%

\bibitem[Sl{\u{a}}tineanu et~al.(2020)]{slatineanu2020}
Sl{\u{a}}tineanu, L., Dodun, O., Cotea{\c{t}}{\u{a}}, M., Nag{\^\i}{\c{t}}, G., B{\u{a}}ncescu, I.\,B., Hri{\c{t}}uc, A., 2020.
Wire electrical discharge machining---a review.
\textit{Machines} 8(4), 69.

\bibitem[Ho et~al.(2004)]{ho2004state}
Ho, K.\,H., Newman, S.\,T., Rahimifard, S., Allen, R.\,D., 2004.
State of the art in wire electrical discharge machining (WEDM).
\textit{International Journal of Machine Tools and Manufacture} 44(12--13), 1247--1259.

\bibitem[Wang et~al.(2023)]{wang2023}
Wang, J., S{\'a}nchez, J.\,A., Izquierdo, B., Ayesta, I., 2023.
Effect of discharge accumulation on wire breakage in WEDM process.
\textit{The International Journal of Advanced Manufacturing Technology} 125(3), 1343--1353.

\bibitem[Luo(1999)]{luo1999}
Luo, Y.\,F., 1999.
Rupture failure and mechanical strength of the electrode wire used in wire EDM.
\textit{Journal of Materials Processing Technology} 94(2--3), 208--215.

\bibitem[Schacht(2004)]{schacht2004}
Schacht, B., 2004.
\textit{Composite Wires and Alternative Dielectrics for Wire Electrical Discharge Machining}.
PhD thesis, Katholieke Universiteit Leuven.

\bibitem[Di~Campli et~al.(2020)]{di_campli_real-time_2020}
Di~Campli, R., Maradia, U., Boccadoro, M., D'Amario, R., Mazzolini, L., 2020.
Real-time Wire EDM Tool Simulation Enabled by Discharge Location Tracker.
\textit{Procedia CIRP} 95, 308--312.

\bibitem[D'Amario(2019)]{damario2019patent}
D'Amario, R., 2019.
Wire electrical discharge machining method.
European Patent EP3446820A1, Agie Charmilles SA.
Available at: \url{https://patents.google.com/patent/EP3446820A1/}.

\bibitem[Gonzalez-Sanchez and Wegener(2025)]{sparc}
Gonzalez-Sanchez, E., Wegener, K., 2025.
SPARC---A Simulation Platform for Advanced Rough-cut Control in Wire EDM.
\textit{Procedia CIRP} 137, 455--461.

\bibitem[Humphreys et~al.(2017)]{humphreys2017}
Humphreys, F.\,J., Rollett, A.\,D., Rohrer, G.\,S., 2017.
\textit{Recrystallization and Related Annealing Phenomena}, 3rd ed.
Elsevier, Oxford.

\bibitem[Zhang(2010)]{zhang2010creep}
Zhang, J.-S., 2010.
Creep Damage Mechanics.
In: \textit{High Temperature Deformation and Fracture of Materials}.
Woodhead Publishing, 274--284.

\bibitem[Larson and Miller(1952)]{larsonmiller1952}
Larson, F.\,R., Miller, J., 1952.
A time-temperature relationship for rupture and creep stresses.
\textit{Transactions of the ASME} 74(5), 765--771.

\bibitem[Stigh(2006)]{stigh2006}
Stigh, U., 2006.
Continuum Damage Mechanics and the Life-Fraction Rule.
\textit{Journal of Applied Mechanics} 73(4), 702--704.

\bibitem[Robinson(1952)]{robinson1952}
Robinson, E.\,L., 1952.
Effect of temperature variation on the long-time rupture strength of steels.
\textit{Transactions of the ASME} 74(5), 777--781.

\bibitem[Kachanov(1986)]{kachanov}
Kachanov, L.\,M., 1986.
\textit{Introduction to Continuum Damage Mechanics}.
Springer, Dordrecht.

\bibitem[ASM(2008)]{asm2008}
ASM, 2008.
\textit{Specialty Handbook: Copper and Copper Alloys}.
ASM International, USA, 35--39.

\bibitem[Ho and Kwok(1989)]{ho1989}
Ho, P.\,S., Kwok, T., 1989.
Electromigration in metals.
\textit{Reports on Progress in Physics} 52(3), 301--348.

\bibitem[Wang and Filippi(2001)]{wang2001}
Wang, P.-C., Filippi, R.\,G., 2001.
Electromigration threshold in copper interconnects.
\textit{Applied Physics Letters} 78(23), 3598--3600.


\bibitem[Pernis et al.(2010)]{pernis2010} Pernis, R., Kasala, J., Bořuta, J., 2010. High temperature plastic deformation of CuZn30 brass calculation of the activation energy. 48(1), 41--46.


\bibitem[Buršıková et al.(2002)]{bursikova2002} Buršıková, V., Svoboda, M., Pahutová, M., Kuchařová, K., Langdon, T.,G., 2002. Creep behaviour of leaded brass. \textit{Materials Science and Engineering: A} 324(1-2), 235--238.


\end{thebibliography}

\clearpage\onecolumn


\end{document}
